\documentclass[11pt,a4paper]{article}

\usepackage[utf8]{inputenc}
\usepackage[T1]{fontenc}
\usepackage{lmodern}
\usepackage{CJKutf8}
\usepackage{microtype}
\usepackage{amsmath, amssymb}
\usepackage{mathtools}
\usepackage{booktabs}
\usepackage{array}
\usepackage{tabularx}
\usepackage{multirow}
\usepackage{xcolor}
\usepackage{hyperref}
\usepackage{geometry}
\usepackage{listings}
\usepackage{enumitem}
\usepackage{titlesec}
\usepackage{abstract}

\geometry{
  top=2.5cm, bottom=2.5cm,
  left=2.5cm, right=2.5cm
}

\hypersetup{
  colorlinks=true,
  linkcolor=blue!60!black,
  citecolor=green!50!black,
  urlcolor=blue!70!black
}

\lstset{
  basicstyle=\small\ttfamily,
  backgroundcolor=\color{gray!10},
  frame=single,
  framesep=4pt,
  breaklines=true,
  columns=fullflexible,
  keepspaces=true,
  xleftmargin=0.5em,
  xrightmargin=0.5em
}

\newcommand{\R}{\mathbb{R}}
\newcommand{\St}{S_t}
\newcommand{\Stp}{S_{t+1}}

\title{
  \textbf{Structural Symmetry Breaking Without Functional Coupling:\\
  Multiplicative FiLM Modulation in Frozen LLM Residual Streams}
}

\author{
  Xiao Ouyang\textsuperscript{1}\thanks{Corresponding author}
  \qquad
  Jing Zhang\textsuperscript{1}
  \\[0.3em]
  \textsuperscript{1}\textit{Wuhan Hongshan Senior High School}
}

\date{\today}

\begin{document}

\maketitle
\thispagestyle{empty}

\begin{abstract}
We investigate how latent state modulation can factorize the residual stream of a
frozen LLM (Qwen3-4B) into a structured hierarchical computation. Adopting
\textbf{multiplicative FiLM} (per-layer feature amplitude gating:
$h_s \times (1 + \alpha\cdot\tanh(W_\ell \cdot S_t))$), we observe a spontaneous
symmetry-breaking transition: inter-layer modulation variance ($w_{\mathrm{std}}$)
grows monotonically from $0.06$ to $0.40$ over 2,650 training steps ($6.6\times$),
far exceeding the additive bottleneck plateau ($w_{\mathrm{std}} = 0.033$).

However, a controlled causal ablation with verified weights reveals a
\textbf{dissociation between structural differentiation and functional influence}:
removing FiLM modulation increases loss by only $\Delta = 0.044$ (pure LM
cross-entropy), while freezing the latent state $\St$ yields $\Delta = 0.0002$. The
dominant source of task improvement ($3.53 \to 2.70$, $-24\%$) comes from
concurrently trained PEFT adapters, not from the modulation mechanism itself.

We conclude that multiplicative FiLM induces genuine layer-wise structural
specialization ($w_{\mathrm{std}}$ growth, gradient decoupling), but this
specialization does not translate into meaningful functional control: layers
differentiate in weight space while remaining functionally near-identical. The
persistent latent state $\St$---the core design premise of the PLAA
architecture---exerts zero measurable influence on output. These findings establish a
critical distinction between \textbf{structural symmetry breaking} and
\textbf{functional causal control} in modulated LLM residual streams, and identify
the effective amplification of learned modulation patterns as a key open problem.
\end{abstract}

\newpage
\setcounter{page}{1}

%------------------------------------------------------------------
\section{Introduction}
%------------------------------------------------------------------

Large language models process all inputs through a uniform residual stream, where
each token re-initializes the computational state. This design achieves remarkable
local reasoning but lacks \textbf{persistent internal state}---there is no mechanism
for maintaining latent variables that evolve across tokens and modulate computation
accordingly.

Previous attempts to inject persistent latent states into LLMs have largely relied on
\textbf{additive gated cross-attention} \cite{sukhbaatar2019augmenting,
wu2022memorizing}, where a learned hidden state $\St$ produces an attention-based
feature added to the residual stream. A companion paper \cite{plaa1} demonstrates
that this additive approach suffers from a fundamental weakness: transformer residual
streams are inherently robust to additive perturbations, absorbing them into
downstream compensation ($\Delta\mathrm{PPL} = 2.01$, below the noise floor of
$\pm 19$).

This paper introduces \textbf{multiplicative FiLM modulation}, where the latent
state controls the \emph{amplitude} of existing residual features via per-layer gain.
We show that this architectural change induces a spontaneous symmetry-breaking
transition in layer-wise modulation structure. However, a carefully controlled causal
ablation---conducted after discovering and correcting a checkpoint persistence
bug---reveals that this structural differentiation is largely \textbf{decoupled from
functional influence} on the model's output.

We initially hypothesized that structural differentiation (measured by
$w_{\mathrm{std}}$) would correlate with functional specialization. Our controlled
ablation disproves this hypothesis, demonstrating that weight-space metrics can be
misleading indicators of computational influence without accompanying causal
verification.

\textbf{Contributions:}

\begin{enumerate}[leftmargin=*, label=\textbf{\arabic*.}]

  \item \textbf{Structural differentiation without functional coupling.}
    Multiplicative FiLM induces $6.6\times$ growth in inter-layer modulation variance
    ($w_{\mathrm{std}}: 0.06 \to 0.40$), but controlled ablation shows causal
    $\Delta = 0.044$---structural symmetry breaking does not entail functional
    specialization.

  \item \textbf{Dissociation of structure and function.}
    Layer-wise weight norms diverge (cosine similarity $0.99 \to 0.34$) while the
    layers' actual computational effect on output remains near-identical. Weight-space
    differentiation and function-space differentiation are separable phenomena.

  \item \textbf{Phase dynamics characterization.}
    We document the monotonic $w_{\mathrm{std}}$ growth trajectory, regime-hopping
    dynamics (intermittent cosine bursts to $0.70$ with self-recovery), and metastable
    coupled hierarchy behavior---independent of the causal coupling magnitude.

  \item \textbf{Negative result on persistent latent state.}
    Freezing the latent state $\St$ changes output by $\Delta = 0.0002$. The
    persistent state---the core design premise of the PLAA architecture---exerts zero
    measurable influence on the model's behavior.

  \item \textbf{Methodological contribution.}
    We identify a critical failure mode in PEFT-based research: \texttt{save\_pretrained()}
    does not persist custom injection layers, leading to silent evaluation against
    zero-initialized weights. We document the corrected evaluation protocol.

\end{enumerate}

%------------------------------------------------------------------
\section{Background and Motivation}
%------------------------------------------------------------------

\subsection{The Additive Bottleneck}

For reference, we summarize the key failure mode of additive injection (studied in
detail in \cite{plaa1}). Gated cross-attention adds a latent-derived signal to the
residual stream:
$$h_s \leftarrow h_s + \mathrm{gate} \cdot e_o, \quad e_o = \mathrm{CrossAttn}(Q{=}h_s,\; K,V{=}\mathrm{proj}(\St))$$
Gate ablation reveals that the maximum causal effect is $\Delta\mathrm{PPL} = 2.01$
at gate scale $2.0$, indistinguishable from additive noise of $\mathrm{std} = \pm 19$.
Multi-stage attenuation is the cause: $\|\St\| \approx 0.15$ (from norm penalty)
yields a tiny projection; $\tanh(0.1) \approx 0.1$ further attenuates; and
$\times 13$ residual stack integrations reduce effective influence to roughly
$0.0003 \cdot h_{\mathrm{hidden}}$.

\subsection{AdaLN FiLM: Preliminary Results}

AdaLN-style modulation \cite{perez2018film, peebles2023scalable} showed apparent
strong causal control in preliminary testing ($\Delta = 11.44$ upon removal).
However, this measurement was conducted under composite loss (including coherence,
repulsion, and norm auxiliary terms) and prior to the discovery of a checkpoint
persistence bug (Section~\ref{sec:checkpoint_bug}). The pure LM causal effect of
AdaLN FiLM remains unverified.

What is established: freezing the state $\St$ barely changes output ($\Delta = 1.33$
under composite loss), confirming that modulation is \textbf{static} regardless of
architecture variant. Attempting to enforce state-dependent modulation via a coupling
loss collapses layer variance ($w_{\mathrm{std}} \to 0$) even at $1/10$ strength,
revealing a fundamental tension between temporal coherence and layer specialization.

%------------------------------------------------------------------
\section{Method}
%------------------------------------------------------------------

\subsection{Architecture}

We inject modulation into layers $16$--$28$ of Qwen3-4B (frozen, 4-bit NF4
quantization, SDPA attention). The latent state $\St \in \R^{256}$ evolves via a GRU
cell conditioned on the final hidden state of each forward pass.

\textbf{Additive baseline (AdaLN FiLM):}
\begin{lstlisting}
hs = LN(hs) * (1 + alpha * tanh(W_scale * S_t))
   + alpha * tanh(W_shift * S_t)
\end{lstlisting}

\textbf{Multiplicative FiLM (this work):}
\begin{lstlisting}
hs = hs * (1 + alpha * tanh(W_l * S_t))
\end{lstlisting}

Key design choices:
\begin{itemize}
  \item \textbf{Per-layer projections.} Each injected layer $\ell$ has its own
    $W_\ell \in \R^{2560 \times 256}$, enabling independent modulation patterns.
  \item \textbf{Zero initialization.} $W_\ell$ initialized to zero, ensuring
    modulation starts as identity ($h_s \times 1$) and specialization must be learned.
  \item \textbf{Multiplicative gating.} The latent state modulates feature amplitude,
    not additive bias.
  \item $\alpha = 0.1$, fixed throughout training.
\end{itemize}

\subsection{State Pipeline}

$$S_0 \;\to\; \text{each turn:}\quad h_s \times (1 + \alpha\cdot\tanh(W_\ell \cdot \St)) \;\to\; \mathrm{forward} \;\to\; \Stp = \mathrm{GRU}(\St,\; h_{\mathrm{last}})$$

The complete layer computation:

\begin{lstlisting}
Input: hidden state h in R^{B x T x 2560}
       latent state S in R^{B x 256}

For each injected layer l in {16...28}:
    h_res  = self_attention(h)          # frozen Qwen3-4B layer
    scale  = tanh(W_l . S)              # per-layer projection
    h_out  = h_res * (1 + 0.1 * scale)  # multiplicative modulation
    h_out  = MLP(h_out)                 # frozen MLP

State update (once per forward pass):
    h_last = h_out[last non-padded token]
    S_{t+1} = GRU(S_t, h_last)
\end{lstlisting}

\subsection{Training Setup}

\textbf{Data.} 3,197 dialogue micro-trajectories (3--5 turns each), stratified as
60\% PersonaChat (neutral continuity), 30\% ESConv (emotional support), 10\%
synthetic roleplay (dynamic range).

\textbf{Loss function.}
$$\mathcal{L} = \mathcal{W}_{\mathrm{LM}}\cdot\mathcal{L}_{\mathrm{LM}} + \mathcal{W}_{\mathrm{COH}}\cdot\|\Delta S\| + \mathcal{W}_{\mathrm{DIV}}\cdot\mathcal{L}_{\mathrm{repulsion}} + \mathcal{W}_{\mathrm{NORM}}\cdot\|S\| + \mathcal{W}_{\mathrm{DIV\_LAYER}}\cdot\mathrm{Var}(\|W_\ell\|)$$

\textbf{Hyperparameters.}

\begin{table}[h]
\centering
\begin{tabular}{ll}
\toprule
\textbf{Hyperparameter} & \textbf{Value} \\
\midrule
Base model         & Qwen3-4B \\
Quantization       & NF4 (4-bit) \\
Attention          & SDPA \\
Frozen layers      & All \\
Trainable params   & 17.5M \\
Optimizer          & AdamW \\
Learning rate      & $1 \times 10^{-4}$ \\
Batch size         & 2 \\
Total steps        & 2,650 \\
$\alpha$ (modulation bound) & 0.1 \\
$\mathcal{W}_{\mathrm{LM}}$ & 0.3 \\
$\mathcal{W}_{\mathrm{COH}}$ & 0.5 \\
$\mathcal{W}_{\mathrm{DIV}}$ & 0.3 \\
$\mathcal{W}_{\mathrm{NORM}}$ & 0.3 \\
$\mathcal{W}_{\mathrm{DIV\_LAYER}}$ & $-0.001$ \\
\bottomrule
\end{tabular}
\caption{Training hyperparameters. $\mathcal{W}_{\mathrm{DIV\_LAYER}}$ is identical
across all architecture variants for controlled comparison.}
\end{table}

\subsection{Metrics}

We define three observables for characterizing the dynamical regime:

\begin{itemize}
  \item $w_{\mathrm{std}}$: Standard deviation of $\|W_\ell\|$ across layers
    $16$--$28$. Measures structural differentiation.
  \item \textbf{cos}: Mean pairwise cosine similarity of gradient updates across
    layers. Measures inter-layer coupling strength.
  \item \textbf{Loss}: Next-token prediction cross-entropy (pure LM loss, excluding
    auxiliary terms). Measures task efficiency.
\end{itemize}

%------------------------------------------------------------------
\section{Results}
%------------------------------------------------------------------

\subsection{The Additive Bottleneck}

We first establish that additive gated cross-attention is causally insufficient.

\begin{table}[h]
\centering
\begin{tabular}{lcc}
\toprule
\textbf{Ablation Condition} & \textbf{PPL} & $\boldsymbol{\Delta}$\textbf{PPL} \\
\midrule
gate $= 0$      & 27.03 & ---  \\
gate $= 1.0$    & 26.26 & 0.77 \\
gate $= 2.0$    & 25.02 & 2.01 \\
\bottomrule
\end{tabular}
\caption{Gate ablation results. Maximum $\Delta = 2.01$, noise
$\mathrm{std} = \pm 19$. Causal influence is indistinguishable from noise.}
\end{table}

\subsection{Multiplicative FiLM Induces Structural Symmetry Breaking}

Switching to multiplicative FiLM triggers an immediate and sustained growth in
layer-wise structural differentiation.

\subsubsection*{Layer Specialization Growth ($w_{\mathrm{std}}$)}

\begin{table}[h]
\centering
\begin{tabular}{cc}
\toprule
\textbf{Step} & $w_{\mathrm{std}}$ \\
\midrule
   50 & 0.062 \\
  350 & 0.244 \\
  800 & 0.333 \\
 1800 & 0.356 \\
 2650 & 0.404 \\
\bottomrule
\end{tabular}
\caption{Monotonic growth of inter-layer modulation variance. $6.6\times$ increase
over 2,650 steps.}
\end{table}

Growth is monotonic across all 53 recorded checkpoints. No saturation plateau is
observed at step 2,650.

\subsubsection*{Phase Plane Dynamics ($w_{\mathrm{std}}$ vs.\ cos)}

The system traces a characteristic trajectory through the
$w_{\mathrm{std}}$--cos phase plane:

\begin{enumerate}
  \item \textbf{Early} (steps 50--300): $w_{\mathrm{std}}$ rises rapidly to $0.24$;
    cos oscillates $0.35$--$0.61$.
  \item \textbf{Mid} (steps 300--1500): $w_{\mathrm{std}}$ plateaus temporarily
    around $0.33$; cos explores $0.34$--$0.63$.
  \item \textbf{Late} (steps 1500--2650): $w_{\mathrm{std}}$ resumes growth to
    $0.40$; cos settles into $0.42$--$0.63$ oscillation.
\end{enumerate}

\subsubsection*{Regime-Hopping Dynamics}

We observe two synchronization burst events (step 1750: cos $= 0.627$; step 1950:
cos $= 0.696$) where global coupling briefly dominates. Three critical observations:

\begin{itemize}
  \item $w_{\mathrm{std}}$ does \emph{not} collapse during bursts
    ($0.348 \to 0.382$, increasing).
  \item Loss recovers to a \emph{lower} baseline after each burst
    (post-step-1950 loss: $10.4 \to 5.3$).
  \item The system's structural differentiation strengthens under perturbation.
\end{itemize}

\subsection{Structure--Performance Correlation}

The $w_{\mathrm{std}}$--loss phase plane reveals a consistent negative correlation
over training: as layer specialization increases, task loss decreases. However,
causal ablation (Section~\ref{sec:causal}) demonstrates that this correlation is
primarily driven by concurrent PEFT adapter learning rather than the modulation
mechanism itself. The FiLM modulation contributes only $\Delta = 0.044$ to task
performance, while the overall improvement ($3.53 \to 2.70$) is attributable to the
LoRA adapters. The $w_{\mathrm{std}}$--loss correlation is therefore
\textbf{spurious with respect to the modulation mechanism}: both metrics improve
over training steps, but without a causal link between them.

\subsection{Causal Ablation with Verified Weights}
\label{sec:causal}

\subsubsection*{The Checkpoint Persistence Bug}
\label{sec:checkpoint_bug}

A critical methodological finding: PEFT's \texttt{save\_pretrained()} only persists
LoRA adapter weights, not custom injection layers. All prior causal ablation
experiments loaded checkpoints with zero-valued \texttt{scale\_proj} weights---the
FiLM modulation was effectively disabled during evaluation without the experimenters'
knowledge.

We corrected this by saving FiLM weights separately:
\begin{lstlisting}
torch.save({
    "plaa_core": plaa_core.state_dict(),
    "scale_proj": {i: layer_blocks[i].scale_proj.state_dict()
                   for i in range(16, 29)},
}, f"{ckpt_dir}/plaa_core.pt")
\end{lstlisting}

Verified checkpoint: scale norms sum $= 74.41$ at step 850,
$w_{\mathrm{std}} = 0.307$. Weights are confirmed non-zero and structurally
differentiated.

\subsubsection*{Ablation Results}

\begin{table}[h]
\centering
\begin{tabular}{lcc}
\toprule
\textbf{Condition} & \textbf{Pure LM Loss} & $\boldsymbol{\Delta}$ \textbf{vs.\ trained} \\
\midrule
Vanilla Qwen3-4B (no adapters, no FiLM) & 3.534 & --- \\
Trained system (LoRA + mFiLM, normal $S$) & 2.700 & $-0.834$ vs.\ vanilla \\
FiLM removed (scale forced to 0) & 2.744 & $+0.044$ \\
State frozen ($S = S_0$, FiLM active) & 2.700 & $+0.0002$ \\
\bottomrule
\end{tabular}
\caption{Causal ablation with verified non-zero FiLM weights. Pure LM
cross-entropy loss, excluding all auxiliary terms.}
\end{table}

\textbf{Three conclusions:}

\begin{enumerate}
  \item \textbf{Overall training is effective.} The system improves over vanilla
    Qwen3-4B by 24\% ($3.53 \to 2.70$), but this improvement is attributable to
    LoRA adapters.
  \item \textbf{FiLM causal effect exists but is minimal.} Removing modulation
    costs only $\Delta = 0.044$. The structural differentiation
    ($w_{\mathrm{std}} = 0.40$) does not translate into functional influence.
  \item \textbf{State coupling is zero.} Freezing $\St$ changes nothing
    ($\Delta = 0.0002$). The persistent latent state is computationally inert.
\end{enumerate}

\subsection{Architecture Comparison}

\begin{table}[h]
\centering
\begin{tabular}{lcccc}
\toprule
\textbf{Architecture} & $w_{\mathrm{std}}$ max & Pure LM loss & cos range & Causal $\Delta$ \\
\midrule
Additive (gated CA)    & 0.033 & ---  & $\sim$0.99   & 2.01 (vs.\ noise)  \\
AdaLN FiLM             & 0.033 & ---  & 0.11--0.24   & 11.44$^\dagger$ \\
\textbf{Mult.\ FiLM (ours)} & \textbf{0.404} & \textbf{2.70} & \textbf{0.34--0.70} & \textbf{0.044} \\
Vanilla Qwen3-4B       & ---   & 3.53 & ---          & --- \\
\bottomrule
\end{tabular}
\caption{Architecture comparison. $^\dagger$Measured under composite loss with
auxiliary terms; pure LM causal effect unverified due to checkpoint persistence bug.}
\end{table}

%------------------------------------------------------------------
\section{Analysis}
%------------------------------------------------------------------

\subsection{Phase Taxonomy}

We propose a four-phase taxonomy based on joint $w_{\mathrm{std}}$--cos--loss
dynamics:

\begin{table}[h]
\centering
\begin{tabular}{clccc p{4.5cm}}
\toprule
\textbf{Phase} & \textbf{Name} & $w_{\mathrm{std}}$ & \textbf{cos} & \textbf{Loss} & \textbf{Dynamics} \\
\midrule
I   & Symmetric             & $<0.05$      & $>0.9$     & High   & Homogeneous layers \\
II  & Convergent            & $0.05$--$0.3$& $0.3$--$0.7$& $\downarrow$ & Emerging specialization \\
III & Specialized           & $0.3$--$0.5$ & $0.3$--$0.5$& Low    & Stable hierarchy \\
IV  & Metastable Coupled    & $0.35$--$0.5$& $0.4$--$0.7$& Low$\pm$ & Hierarchy + bursts \\
\bottomrule
\end{tabular}
\caption{Four-phase taxonomy. Our system occupies Phase IV: a metastable coupled
hierarchy. Note that phase classification is based on structural metrics only;
causal ablation reveals that Phase IV structural differentiation does not imply
functional specialization.}
\end{table}

\subsection{Why Structure Diverges from Function}

The dissociation between $w_{\mathrm{std}}$ growth and causal $\Delta$ admits a
clear mechanistic explanation. With $\alpha = 0.1$ and $\tanh$ bounding, the
effective modulation range is $\pm 10\%$ of feature amplitude. Even with maximally
differentiated layer weights, the perturbation to the residual stream is:
$$\|h_s \cdot 0.1 \cdot \tanh(W_\ell \cdot S)\| \leq 0.1 \cdot \|h_s\|$$
This is small relative to the residual stream magnitude ($\|h_s\| \sim 50$--$100$),
and downstream layers can compensate for perturbations of this scale. The
$w_{\mathrm{DIV\_LAYER}}$ loss encourages weight differentiation, but the bounded
modulation prevents this differentiation from producing functionally distinct
computations.

\subsection{Metastability as a Structural (Not Functional) Phenomenon}

The regime-hopping dynamics---synchronization bursts followed by recovery to
lower-loss attractors---are genuine dynamical phenomena in the training process.
However, given that FiLM's causal contribution is only $\Delta = 0.044$, these
dynamics primarily reflect the \textbf{LoRA adapter's} optimization landscape rather
than the modulation mechanism's. The FiLM layers are structurally evolving but
functionally passive passengers in the training process.

\subsection{Implications for Latent Modulation Research}

Our findings carry a methodological warning: \textbf{weight-space metrics
($w_{\mathrm{std}}$, gradient cosine) can be misleading indicators of functional
influence}. A modulation mechanism can exhibit rich structural dynamics---symmetry
breaking, phase transitions, regime hopping---while contributing negligibly to the
model's actual computation. Causal ablation with verified weights is essential for
any claim of functional influence.

%------------------------------------------------------------------
\section{Related Work}
%------------------------------------------------------------------

\textbf{Conditional feature modulation.} FiLM \cite{perez2018film} and AdaIN
\cite{huang2017arbitrary} introduce multiplicative conditioning in vision networks.
DiT \cite{peebles2023scalable} applies AdaLN to diffusion transformers. Our work
extends multiplicative conditioning to LLM residual streams and provides the first
controlled causal evaluation of its functional effect.

\textbf{Layer specialization in transformers.} Prior work has documented functional
specialization across transformer layers \cite{tenney2019bert, clark2019does,
rogers2020primer}. We show that structural specialization (weight divergence) can
emerge spontaneously through architectural pressure without producing corresponding
functional specialization---a distinction not previously documented.

\textbf{Latent state injection.} Sukhbaatar et al.\ \cite{sukhbaatar2019augmenting}
and Wu et al.\ \cite{wu2022memorizing} inject persistent state additively. We
demonstrate that both additive and multiplicative injection fail to achieve
meaningful state-driven control under bounded modulation.

\textbf{Non-equilibrium dynamics in neural networks.} Gur-Ari et al.\
\cite{gurari2018gradient} characterize gradient dynamics of wide networks. Our
regime-hopping observations extend this to modulation layers, though we note that
these dynamics are structural rather than functionally consequential.

%------------------------------------------------------------------
\section{Conclusion}
%------------------------------------------------------------------

We demonstrate that multiplicative FiLM modulation induces spontaneous symmetry
breaking in frozen LLM residual streams, driving layer-wise structural
differentiation ($w_{\mathrm{std}}: 0.06 \to 0.40$, $6.6\times$ growth) with
metastable coupled hierarchy dynamics. However, controlled causal ablation with
verified weights reveals that this structural differentiation is \textbf{decoupled
from functional influence}: removing FiLM costs only $\Delta = 0.044$, and freezing
the latent state $\St$ yields $\Delta = 0.0002$.

\textbf{Key finding.} Weight-space symmetry breaking and output-space causal control
are separable phenomena. A modulation mechanism can induce genuine structural
specialization (layers learn different weight patterns) while contributing negligibly
to the model's actual computation. Metrics like $w_{\mathrm{std}}$, which track
weight-space differentiation, can be misleading indicators of functional influence
without accompanying causal ablation.

\textbf{Limitations.} The $\alpha = 0.1$ tanh bound constrains effective modulation
to $\pm 10\%$ of feature amplitude. This may be insufficient for the modulation to
overcome the residual stream's inherent robustness. The persistent latent state
$\St$---the core premise of the PLAA architecture---is computationally inert under
current training.

\textbf{Future work.} Three directions may bridge the structure--function gap:
(1)~increasing $\alpha$ or removing the $\tanh$ bound to amplify modulation
magnitude; (2)~gated routing mechanisms that force the model to \emph{depend} on
modulation signals for correct prediction; (3)~competitive computation paradigms
\cite{shazeer2017outrageously} where modulation controls which expert processes each
token. The central open problem is clear: how to convert structural differentiation
into functional specialization.

%------------------------------------------------------------------


\appendix

\section{Complete Training Data}
All 53 data points at 50-step intervals (steps 50 to 2650) are provided in supplementary materials.

\section{Architecture Detail}
See paper body for full architecture specification.

\begin{thebibliography}{99}
\bibitem{perez2018film} Perez, E., et al. FiLM: Visual Reasoning with a General Conditioning Layer. AAAI 2018.
\bibitem{peebles2023scalable} Peebles, W., & Xie, S. Scalable Diffusion Models with Transformers. ICCV 2023.
\bibitem{sukhbaatar2019augmenting} Sukhbaatar, S., et al. Augmenting Self-attention with Persistent Memory. arXiv:1907.01470.
\bibitem{wu2022memorizing} Wu, Y., et al. Memorizing Transformers. ICLR 2022.
\bibitem{huang2017arbitrary} Huang, X., & Belongie, S. Arbitrary Style Transfer in Real-time with Adaptive Instance Normalization. ICCV 2017.
\bibitem{tenney2019bert} Tenney, I., et al. BERT Rediscovers the Classical NLP Pipeline. ACL 2019.
\bibitem{clark2019does} Clark, K., et al. What Does BERT Look At? BlackboxNLP 2019.
\bibitem{rogers2020primer} Rogers, A., et al. A Primer in BERTology. TACL 2020.
\bibitem{gurari2018gradient} Gur-Ari, G., et al. Gradient Descent Happens in a Tiny Subspace. arXiv:1812.04754.
\bibitem{shazeer2017outrageously} Shazeer, N., et al. Outrageously Large Neural Networks: The Sparsely-Gated Mixture-of-Experts Layer. ICLR 2017.
\bibitem{plaa1} Ouyang, X. & Zhang, J. PLAA: Persistent Latent Affective Architecture. Companion paper.
\end{thebibliography}

\end{document}
